\section{Cómo el sonido continuo se convierte en «lenguaje»: VQ, Codebook y Hierarchy}

El segundo trabajo pasa de la generación a la representation learning. El obstáculo central es el siguiente: un modelo de lenguaje predice sobre un vocabulario finito, mientras que la waveform/embedding de audio es continua y de alta frecuencia. Si se quiere usar un next-token objective, primero hay que definir el «audio token». La Vector quantization asigna el encoder output continuo a un index discreto mediante un codebook, lo que hace posible el sequence modeling, pero al mismo tiempo introduce problemas nuevos como el reconstruction error, la code utilization y el collapse.

\lecturefigure{slide-25-contrastive-audio-representations-title.jpg}{Segundo trabajo: generative y contrastive audio representation learning.}{Video oficial de Stanford Online 00:22:36--00:22:44; Verma y Smith, arXiv:2010.11459.}

\subsection{Por qué aprender una Audio Representation}

La Acoustic scene understanding busca identificar el entorno o los eventos a partir del sonido, sin necesidad de generar cada sample. El ponente da ejemplos como emergency vehicle, music recommendation y health monitoring: una buena representation debe comprimir la estructura relevante para la decisión y mantener una invariancia adecuada frente a los detalles de grabación irrelevantes. Sin embargo, estas aplicaciones también tienen costos distintos de falsos positivos, por lo que no basta con usar una accuracy uniforme en lugar de la calidad real de las decisiones.

\lecturefigure{slide-26-applications-overview.jpg}{Aplicaciones de la Acoustic scene understanding: tránsito, recomendación, medicina y percepción del entorno.}{Video oficial de Stanford Online 00:22:45--00:24:32.}

\teachervoice{El ejemplo de los micrófonos de Waymo subraya que el sonido puede revelar una sirena con anticipación, fuera del campo visual; el ejemplo de la tos y los estornudos en medicina recuerda que la representation llega a decisiones de alto riesgo. El significado aguas abajo determina qué información vale la pena conservar.}

\subsection{Del espacio continuo a un vocabulario finito}

Sea un encoder que asigna el fragmento $x$ a $z_e(x)\in\mathbb{R}^d$, con codebook $\{e_1,\ldots,e_K\}$. La Vector quantization elige el code más cercano:
\[
k^*(x)=\operatorname*{argmin}_{k\in\{1,\ldots,K\}}
\|z_e(x)-e_k\|_2^2,
\qquad z_q(x)=e_{k^*(x)}.
\]
El token discreto es el index $k^*$, y el decoder reconstruye la waveform o el spectrum a partir de $z_q$. Lo que ve el Transformer es la index sequence, no el sonido original; el paso temporal y la capacidad del codebook determinan el token rate y el límite superior de la reconstrucción.

\lecturefigure{slide-27-discrete-audio-language-modeling.jpg}{Requisito previo del Audio language modeling: convertir primero el sonido continuo en una code sequence discreta.}{Video oficial de Stanford Online 00:24:33--00:25:57.}

\lecturefigure{slide-28-vector-quantization-recipe.jpg}{VQ/VQ-VAE, k-means y los métodos tipo wav2vec ofrecen distintas vías de discretización.}{Video oficial de Stanford Online 00:25:58--00:26:53.}

\begin{knowledgebox}{Términos clave: los cuatro objetos de un sistema de Codebook}
El encoder determina el feature continuo; el quantizer elige el code más cercano; el codebook almacena un número finito de prototypes; el decoder determina si la información conservada en el code permite reconstruir la señal. Un modelo de lenguaje puede predecir el Code index, pero su «semántica» depende del objetivo de entrenamiento y no tiene por qué corresponder a fonemas, notas musicales ni eventos con nombre humano.
\end{knowledgebox}

\subsection{Voronoi cell, reconstrucción y Codebook Collapse}

Desde el punto de vista geométrico, cada code vector posee una Voronoi region, y los encoder outputs que caen en esa región comparten el mismo index. Un $K$ demasiado pequeño produce un quantization error alto; un $K$ demasiado grande o una optimización desequilibrada hace que muchos codes nunca se elijan, lo que da lugar al codebook collapse. El entrenamiento también debe equilibrar reconstruction, codebook update y commitment, para que el encoder no oscile arbitrariamente entre las fronteras de los codes.

\lecturefigure{slide-29-vq-voronoi.jpg}{La Vector quantization vista desde Voronoi: el espacio continuo se divide según el code más cercano.}{Video oficial de Stanford Online 00:26:54--00:27:55.}

\lecturefigure{slide-30-vq-autoencoder-transformer.jpg}{Encode--quantize--predict--decode: la combinación de un VQ autoencoder y un Transformer.}{Video oficial de Stanford Online 00:27:56--00:29:55.}

\begin{lstlisting}[caption={Flujo conceptual del modelado con audio-tokens discretos.}]
def audio_token_pipeline(waveform, encoder, codebook,
                         transformer, decoder):
    continuous = encoder(waveform)
    token_ids = nearest_code_indices(continuous, codebook)
    predicted_ids = transformer.autoregressive(token_ids)
    quantized = codebook[predicted_ids]
    reconstructed = decoder(quantized)
    return token_ids, reconstructed
\end{lstlisting}

\begin{warningbox}{Los tokens discretos no son una compresión gratuita}
Un token rate más bajo reduce la longitud que procesa el Transformer, pero pierde timing, phase, timbre o transitorios pequeños; un token rate más alto mejora la reconstrucción, pero vuelve a introducir secuencias largas. Al evaluar un codec/tokenizer hay que reportar a la vez la calidad de reconstrucción, la code utilization, el bitrate y el desempeño aguas abajo.
\end{warningbox}

\subsection{Hierarchical Codes: cada nivel se encarga de una escala temporal distinta}

Los sistemas del tipo Jukebox usan VQ de varios niveles: los tokens de nivel alto y baja frecuencia capturan la estructura musical más lenta, los tokens de nivel bajo y alta frecuencia completan los detalles acústicos locales, y después el decoder sintetiza la waveform. La Hierarchy convierte la tensión multiescala vista antes en varias token streams; el costo es que el entrenamiento, el sampling y la coherencia entre niveles se vuelven más complejos.

\lecturefigure{slide-31-jukebox-hierarchical-tokens.jpg}{Jukebox-style hierarchy: los codes de varios niveles se reparten el trabajo entre la compresión temporal y los detalles de reconstrucción.}{Video oficial de Stanford Online 00:29:56--00:30:24.}

\begin{knowledgebox}{Lectura de la figura: por qué la Hierarchy reduce la long-range burden}
Si cada token del nivel más alto cubre un intervalo de tiempo más largo, el Transformer puede modelar la estructura a nivel de canción con una sequence más corta; los niveles inferiores solo tienen que completar los detalles bajo la condition del nivel superior. No elimina el cómputo, sino que asigna la planeación global y la reconstrucción local a resoluciones distintas.
\end{knowledgebox}

\teachervoice{El ponente califica la VQ como «powerful and under-utilized». Lo que realmente vale la pena conservar de la clase no es el calificativo, sino la forma de combinar: convertir el continuous audio en codes y luego hacer que el Transformer aprenda la long dependency de esos codes.}

\subsection{Resumen de la sección}

El Audio language modeling depende del tokenizer. La VQ permite que el sonido continuo entre en un vocabulario finito, pero el token rate, el codebook size y el decoder determinan en conjunto qué estructura de largo alcance puede aprenderse y qué detalles acústicos pueden recuperarse. La Hierarchy es una asignación explícita de ancho de banda en varias escalas.
